% =============================================================
%  Relatorio de Pesquisa de Satisfacao - PetCuida (versao Alpha 3i)
%  Registro formalizado de entrevista semiestruturada
%  Compilar com: pdflatex entrevista_petcuida_flavia.tex (2x)
% =============================================================
\documentclass[12pt,a4paper]{article}

\usepackage[utf8]{inputenc}
\usepackage[T1]{fontenc}
\usepackage[brazil]{babel}
\usepackage{lmodern}
\usepackage[margin=3cm]{geometry}
\usepackage{setspace}
\usepackage{booktabs}
\usepackage{longtable}
\usepackage{array}
\usepackage{enumitem}
\usepackage{csquotes}
\usepackage{fancyhdr}
\usepackage{titlesec}
\usepackage[hidelinks]{hyperref}

\onehalfspacing

\pagestyle{fancy}
\fancyhf{}
\fancyhead[L]{\footnotesize Pesquisa de satisfação --- PetCuida \textit{Alpha 3i}}
\fancyhead[R]{\footnotesize \thepage}
\renewcommand{\headrulewidth}{0.4pt}

% Ambiente para os turnos de fala da transcricao
\newenvironment{turnos}
  {\begin{list}{}{%
     \setlength{\leftmargin}{2.2cm}%
     \setlength{\labelwidth}{2.0cm}%
     \setlength{\labelsep}{0.2cm}%
     \setlength{\itemsep}{0.6em}%
     \setlength{\parsep}{0pt}}}
  {\end{list}}

\newcommand{\fala}[1]{\item[\textbf{#1}\hfill]}

\newcommand{\entrevistador}{Giovanna}
\newcommand{\participante}{Sra. Flavia}

\title{\textbf{Avaliação de Satisfação e Usabilidade Percebida do\\
Aplicativo PetCuida (versão \textit{Alpha 3i})}\\[0.4cm]
\large Registro formalizado de entrevista semiestruturada com usuária potencial}

\author{Entrevistadora: \entrevistador{} (aluna responsável pela coleta)\\
Participante: \participante{} (usuária potencial)}

\date{Belo Horizonte, 16 de setembro de 2026}

\begin{document}

\maketitle
\thispagestyle{empty}

\begin{abstract}
\noindent Este documento formaliza, para fins acadêmicos, o registro de uma entrevista
semiestruturada realizada com uma usuária potencial do aplicativo \textbf{PetCuida}, em sua
versão \textit{Alpha 3i}. A sessão teve por objetivo apresentar a proposta funcional do
software e coletar a percepção da participante quanto à interface, às funcionalidades
apresentadas, à utilidade percebida e às possibilidades de expansão da solução. Os dados
foram obtidos de forma remota, por meio de chamada de vídeo, após a participante assistir
a uma demonstração gravada do protótipo funcional. A participante apresentou avaliação
positiva da interface, destacando seu caráter atrativo e a utilização de elementos gráficos.
Também foram identificadas percepções favoráveis em relação às funcionalidades de vacinação,
prontuário, calendário, suporte humano, notificações, triagem de sintomas e localização de
clínicas. Como pontos de atenção, a participante relatou não compreender inicialmente o
funcionamento dos créditos e questionou a finalidade do QR Code do pet. Foram ainda
sugeridas novas funcionalidades relacionadas à alimentação durante períodos de doença,
adoção de animais, lojas de produtos para pets, promoções e comparação de preços.

\vspace{0.4cm}
\noindent\textbf{Palavras-chave:} usabilidade percebida; pesquisa de satisfação; entrevista
semiestruturada; aplicativos móveis; bem-estar animal; PetCuida.
\end{abstract}

\newpage
\tableofcontents
\newpage

% =============================================================
\section{Identificação da Coleta}

\begin{table}[h!]
\centering
\begin{tabular}{@{}>{\raggedright\arraybackslash}p{5.2cm}>{\raggedright\arraybackslash}p{8.5cm}@{}}
\toprule
\textbf{Item} & \textbf{Descrição} \\
\midrule
Objeto avaliado & Aplicativo PetCuida - versão \textit{Alpha 3i} (protótipo funcional Flutter/Android) \\
Instrumento & Entrevista semiestruturada com demonstração gravada do protótipo \\
Entrevistadora & \entrevistador{} (aluna responsável) \\
Participante & \participante{} \\
Perfil da participante & Usuária potencial, aproximadamente 50 anos, tutora de um cachorro de pequeno porte \\
Data da coleta & 16 de setembro de 2026 \\
Duração aproximada & Duração do vídeo de demonstração acrescida de aproximadamente 3 minutos de entrevista \\
Local & Entrevista remota realizada por chamada de vídeo \\
Registro & Vídeo de demonstração e registro da entrevista \\
Modalidade de análise & Análise qualitativa de conteúdo, por categorização temática \\
\bottomrule
\end{tabular}
\caption{Ficha técnica da coleta de dados.}
\end{table}

% =============================================================
\section{Objetivo da Pesquisa}

A presente coleta teve como objetivo geral \textbf{verificar a aceitação e a satisfação
percebida de uma usuária potencial diante da proposta funcional do aplicativo PetCuida em
estágio alfa}. Como objetivos específicos, delimitam-se:

\begin{enumerate}[label=\alph*)]
  \item apresentar à participante o conceito do produto e suas principais funcionalidades;
  \item permitir que a participante conhecesse o fluxo de telas implementado na versão
        \textit{Alpha 3i} por meio de uma demonstração gravada;
  \item registrar a percepção da participante quanto à interface e à apresentação visual
        do software;
  \item identificar as funcionalidades consideradas mais úteis ou relevantes;
  \item coletar pontos de atenção, dúvidas e sugestões de melhoria;
  \item identificar novas funcionalidades consideradas desejáveis pela participante.
\end{enumerate}

Ressalta-se que \textbf{esta rodada não constituiu teste de usabilidade com execução de
tarefas pelo usuário}, uma vez que a participante assistiu a uma demonstração previamente
gravada do protótipo e, posteriormente, participou de uma conversa remota para fornecer
suas percepções.

% =============================================================
\section{Descrição do Objeto Avaliado}

O PetCuida é um aplicativo móvel cuja proposta consiste em intermediar, em uma mesma
plataforma, a demanda e a oferta de serviços de cuidado com animais de estimação. O modelo
foi verbalizado por meio da analogia de um \enquote{\textit{Uber} de mão dupla}: o mesmo
usuário pode figurar tanto como solicitante quanto como prestador de serviços - tais como
passeio, hospedagem e transporte do animal -, enquanto profissionais da medicina
veterinária e clínicas conveniadas podem prestar atendimento pela plataforma.

A contrapartida pela prestação de serviços é convertida em \textbf{crédito} utilizável pelo
próprio usuário para custear os cuidados de seu animal, constituindo o eixo de acessibilidade
econômica da proposta.

A versão \textit{Alpha 3i} implementa, com dados locais e sem dependência de
\textit{backend}, o fluxo completo do diagrama de telas, conforme apresentado na
Tabela~\ref{tab:modulos}.

\begin{table}[h!]
\centering
\small
\begin{tabular}{@{}>{\raggedright\arraybackslash}p{4.2cm}>{\raggedright\arraybackslash}p{9.5cm}@{}}
\toprule
\textbf{Módulo} & \textbf{Funcionalidades demonstradas na sessão} \\
\midrule
Autenticação & Tela de carregamento; acesso por usuário e senha; acesso e registro por conta Google; criação de conta. \\
\textit{Onboarding} do pet & Seleção de espécie; cadastro de nome; campo opcional de raça; seleção de data de nascimento; definição de foto por galeria ou captura imediata. \\
Área do pet & Painel de cuidados; situação vacinal; calendário de vacinas; agendamentos; prontuário do animal; edição dos dados cadastrados. \\
Triagem de serviço & Pré-seleção orientativa da necessidade do usuário, com possibilidade de interação conversacional. \\
Rede solidária & Listagem de serviços disponíveis; mapa esquemático de pessoas e prestadores próximos; solicitação de serviço. \\
Perfil e configurações & Dados do tutor; definição de foto; informações do programa; termos de uso; políticas de privacidade. \\
\bottomrule
\end{tabular}
\caption{Módulos do PetCuida \textit{Alpha 3i} percorridos durante a demonstração.}
\label{tab:modulos}
\end{table}

% =============================================================
\section{Procedimentos Metodológicos}

\subsection{Abordagem}

Adotou-se abordagem \textbf{qualitativa}, de caráter \textbf{exploratório}, por meio de
entrevista semiestruturada conduzida remotamente. A opção por esse método decorre do
estágio inicial de maturidade do artefato: em versões alfa, interessa menos a mensuração
estatística de desempenho e mais a captura da compreensão, da expectativa e da reação
espontânea do público-alvo diante da proposta de valor.

\subsection{Instrumento e condução}

A sessão foi organizada em três momentos:

\begin{enumerate}[label=\textbf{\arabic*.}]
  \item \textbf{Apresentação} - a participante assistiu a um vídeo gravado contendo a
        demonstração do aplicativo PetCuida e de suas principais funcionalidades;
  \item \textbf{Observação} - durante a apresentação, a participante acompanhou os fluxos
        e recursos disponíveis na versão \textit{Alpha 3i};
  \item \textbf{Coleta de percepção} - após a demonstração, foi realizada uma conversa
        remota de aproximadamente três minutos, organizada a partir de quatro tópicos:
        percepção geral da interface, funcionalidades de destaque, pontos positivos e
        possibilidades de melhoria, e sugestões de novas funcionalidades.
\end{enumerate}

\subsection{Roteiro aplicado}

O roteiro foi organizado em quatro perguntas abertas:

\begin{enumerate}[label=\textbf{Q\arabic*.},leftmargin=1.6cm]
  \item Após assistir à demonstração do aplicativo, qual é a sua percepção sobre a interface
        e a facilidade de uso?
  \item Quais funcionalidades apresentadas chamaram mais a sua atenção ou você considera
        mais úteis?
  \item Existe alguma funcionalidade que você considera especialmente importante ou algum
        ponto que poderia ser melhorado?
  \item Existe alguma outra funcionalidade que você gostaria de encontrar em um aplicativo
        desse tipo?
\end{enumerate}

\subsection{Tratamento dos dados}

O registro da entrevista foi analisado qualitativamente, considerando as manifestações
espontâneas da participante após a apresentação do protótipo. Para fins de sistematização,
as respostas foram organizadas em quatro tópicos correspondentes ao roteiro da entrevista.

A transcrição formalizada apresentada na Seção~\ref{sec:transcricao} preserva o conteúdo
proposicional das manifestações da participante, realizando apenas edição de superfície
para adequação à redação acadêmica, sem alteração do sentido das respostas.

% =============================================================
\section{Transcrição Formalizada da Entrevista}
\label{sec:transcricao}

\noindent\textit{Convenções:} \textbf{E} designa a entrevistadora (\entrevistador{}) e
\textbf{P} designa a participante (\participante{}). As falas foram organizadas a partir
dos quatro tópicos do roteiro. Intervenções sem conteúdo proposicional foram suprimidas.

\begin{turnos}

\fala{E:} Após assistir à demonstração do aplicativo, qual é a sua percepção sobre a
interface e a facilidade de uso?

\fala{P:} O processo de login achei bem normal, mas a interface é bonita e atrativa,
com desenhos super legais.

\fala{E:} Quais funcionalidades apresentadas chamaram mais a sua atenção ou você considera
mais úteis?

\fala{P:} Achei legal a questão das vacinas, prontuário e calendário. O QR Code do pet
eu não sei bem para que serve, mas vi também aqui. Os créditos eu não entendi bem como
funcionam. Achei o tutor ter foto ali muito legal.

\fala{E:} Existe alguma funcionalidade que você considera especialmente importante ou
algum ponto que poderia ser melhorado?

\fala{P:} Achei legal ter suporte humano, porque é importante. Notificações e lembretes
ajudam muito o dono do pet. Eu acho que deveria ter um controle de alimentação para pets
quando ele fica doente. A triagem de sintomas acho muito legal mesmo, porque aí a gente
pode tentar entender o problema do pet. O mapa de clínicas é muito útil, principalmente
se tiver horário.

\fala{E:} Existe alguma outra funcionalidade que você gostaria de encontrar em um aplicativo
desse tipo?

\fala{P:} Poderia ter algo de adoção também de pets, lista de lojas para comprar artigos
de pet, oferta de promoções destas lojas participantes e pesquisa de preços de artigos de
pet mais baratos no mercado, onde a gente coloca o tipo de produto que quer e ele mostra
todas as lojas que têm para vender e quanto custa.

\end{turnos}

% =============================================================
\section{Análise dos Resultados}

\subsection{Categorização temática}

A partir das falas da participante, identificaram-se seis categorias principais de análise,
sintetizadas na Tabela~\ref{tab:categorias}.

\begin{table}[h!]
\centering
\small
\begin{tabular}{@{}>{\raggedright\arraybackslash}p{3.5cm}>{\raggedright\arraybackslash}p{5.3cm}>{\raggedright\arraybackslash}p{4.5cm}@{}}
\toprule
\textbf{Categoria} & \textbf{Evidência na fala} & \textbf{Interpretação} \\
\midrule

Percepção visual positiva &
\enquote{a interface é bonita e atrativa}; \enquote{desenhos super legais} &
A participante demonstrou percepção positiva em relação à apresentação visual do aplicativo,
associando a interface a uma experiência atrativa. \\

\addlinespace

Gestão da saúde do pet &
\enquote{vacinas, prontuário e calendário}; \enquote{notificações e lembretes} &
As funcionalidades relacionadas ao acompanhamento da saúde e à organização dos cuidados
foram percebidas como relevantes para a rotina do tutor. \\

\addlinespace

Suporte e orientação &
\enquote{suporte humano}; \enquote{triagem de sintomas} &
A participante valorizou recursos que oferecem orientação ao tutor, especialmente diante
de dúvidas ou possíveis problemas de saúde do animal. \\

\addlinespace

Utilidade do mapa de clínicas &
\enquote{O mapa de clínicas é muito útil, principalmente se tiver horário} &
A funcionalidade de localização foi considerada útil, com destaque para a possibilidade
de consultar os horários de atendimento. \\

\addlinespace

Dificuldades de compreensão &
\enquote{não sei bem para que serve}; \enquote{não entendi bem como funcionam} &
O QR Code do pet e o sistema de créditos não foram compreendidos de forma imediata,
indicando necessidade de maior clareza na comunicação dessas funcionalidades. \\

\addlinespace

Ampliação do ecossistema &
\enquote{adoção}; \enquote{lista de lojas}; \enquote{promoções}; \enquote{pesquisa de preços} &
A participante sugeriu ampliar o aplicativo para além dos cuidados veterinários, incorporando
serviços de adoção e recursos relacionados à aquisição de produtos para pets. \\

\bottomrule
\end{tabular}
\caption{Categorias emergentes da análise de conteúdo da entrevista com a participante.}
\label{tab:categorias}
\end{table}

\subsection{Discussão}

A avaliação da participante apresentou percepção positiva em relação à interface do
PetCuida. O destaque dado à aparência visual, caracterizada como \enquote{bonita e
atrativa}, indica que os elementos gráficos utilizados na apresentação do aplicativo
contribuíram para uma primeira impressão favorável.

Entre as funcionalidades existentes, a participante destacou principalmente os recursos
relacionados à gestão da saúde do animal, como vacinação, prontuário e calendário. Também
foram valorizados recursos de notificações e lembretes, considerados úteis para auxiliar
o tutor na organização dos cuidados com o pet.

A presença de suporte humano foi considerada importante pela participante. De maneira
semelhante, a triagem de sintomas foi percebida como uma funcionalidade relevante por
possibilitar uma tentativa inicial de compreensão de possíveis problemas enfrentados pelo
animal. O mapa de clínicas também recebeu avaliação positiva, especialmente quando
associado à possibilidade de apresentar os horários de atendimento dos estabelecimentos.

Por outro lado, foram identificados dois pontos de compreensão que merecem atenção. A
participante declarou não compreender bem a finalidade do QR Code do pet e afirmou não
ter entendido inicialmente como funcionam os créditos. Esses comentários indicam que
determinadas funcionalidades podem necessitar de explicações adicionais ou de uma
apresentação mais clara dentro da interface.

Por fim, a participante apresentou diversas sugestões de expansão da proposta. Entre elas,
destacam-se um controle de alimentação específico para períodos em que o pet esteja doente,
um espaço destinado à adoção de animais, uma lista de lojas de produtos para pets,
promoções de estabelecimentos participantes e uma ferramenta de comparação de preços.
Essas sugestões apontam para uma expectativa de que o aplicativo possa reunir, em uma
mesma plataforma, diferentes necessidades relacionadas à rotina de tutores de animais.

% =============================================================
\subsection{Pontos de atenção identificados}

A partir da entrevista, os principais pontos que podem ser considerados para futuras
iterações do artefato são:

\begin{itemize}[leftmargin=*]
  \item fornecer uma explicação mais clara sobre a finalidade do QR Code associado ao pet;
  \item incluir explicações sobre o funcionamento e a utilização dos créditos;
  \item avaliar a possibilidade de apresentar horários de funcionamento junto ao mapa de
        clínicas;
  \item considerar a inclusão de recursos de controle de alimentação, especialmente em
        situações de doença do animal;
  \item avaliar a criação de uma área destinada à adoção de pets;
  \item avaliar a inclusão de lojas participantes e ofertas promocionais;
  \item estudar a viabilidade de uma funcionalidade de comparação de preços de produtos
        para pets.
\end{itemize}

% =============================================================
\section{Limitações Metodológicas}

O presente registro deve ser lido com as seguintes ressalvas:

\begin{enumerate}[leftmargin=*]
  \item \textbf{Amostra reduzida.} A entrevista representa a percepção de uma única
        participante e, portanto, não permite generalização estatística. Os resultados
        possuem caráter indicativo e exploratório.

  \item \textbf{Avaliação baseada em demonstração.} A participante assistiu a um vídeo
        contendo a demonstração do protótipo e não realizou tarefas diretamente no
        aplicativo. Dessa forma, as observações registradas representam percepção e
        expectativa, e não métricas objetivas de usabilidade.

  \item \textbf{Ausência de instrumento padronizado.} Não foi aplicada escala validada,
        como SUS ou escala Likert de satisfação, o que impede uma comparação quantitativa
        direta com outras versões do software.

  \item \textbf{Tempo reduzido de interação.} A etapa de conversa após a demonstração
        teve duração aproximada de três minutos, restringindo a quantidade de questões
        que poderiam ser aprofundadas.
\end{enumerate}

% =============================================================
\section{Considerações Finais e Recomendações}

A entrevista cumpriu o objetivo de coletar a percepção de uma potencial usuária sobre a
versão \textit{Alpha 3i} do PetCuida. A participante apresentou percepção positiva quanto
à interface visual e identificou como relevantes diversas funcionalidades relacionadas
à saúde e ao acompanhamento do animal.

Os principais pontos de atenção encontrados estão relacionados à compreensão de determinadas
funcionalidades, especialmente o QR Code do pet e o sistema de créditos. Ao mesmo tempo,
a participante apresentou sugestões de expansão que podem contribuir para ampliar o
conjunto de necessidades atendidas pelo aplicativo.

Para as próximas rodadas, recomenda-se:

\begin{enumerate}[label=\textbf{R\arabic*.},leftmargin=1.6cm]
  \item ampliar a amostra, contemplando diferentes perfis de tutores de animais;
  \item realizar testes de usabilidade com execução direta de tarefas no protótipo;
  \item avaliar a inclusão de explicações sobre o sistema de créditos;
  \item tornar mais clara a finalidade do QR Code do pet;
  \item estudar a apresentação de horários de atendimento no mapa de clínicas;
  \item avaliar a viabilidade de recursos de controle de alimentação durante períodos
        de doença do animal;
  \item analisar a possibilidade de incluir uma área de adoção de pets;
  \item avaliar a inclusão de lojas, promoções e mecanismos de comparação de preços de
        produtos para animais.
\end{enumerate}

% =============================================================
\appendix

\section{Termo de Registro e Uso dos Dados}

A participante foi informada sobre a finalidade acadêmica da coleta e sobre o uso restrito
dos dados para fins de avaliação do artefato desenvolvido, manifestando concordância.
Dados pessoais identificáveis foram tratados de forma restrita no registro da pesquisa.

\section{Ficha Técnica do Artefato Avaliado}

\begin{table}[h!]
\centering
\small
\begin{tabular}{@{}>{\raggedright\arraybackslash}p{4.6cm}>{\raggedright\arraybackslash}p{9.0cm}@{}}
\toprule
\textbf{Item} & \textbf{Descrição} \\
\midrule
Denominação & PetCuida \\
Versão avaliada & \textit{Alpha 3i} \\
Plataforma & Android \\
Tecnologia & Flutter/Dart, com navegação por rotas e gerência de estado por notificadores \\
Persistência & Dados locais e demonstrativos, sem \textit{backend} \\
Escopo implementado & Autenticação, \textit{onboarding}, painel inicial, área do pet, triagem, rede solidária e perfil do tutor \\
\bottomrule
\end{tabular}
\caption{Ficha técnica resumida da versão avaliada.}
\end{table}

\section{Fontes Primárias do Registro}

\begin{itemize}[leftmargin=*]
  \item Vídeo de demonstração da versão \textit{Alpha 3i} do aplicativo;
  \item Registro da entrevista remota realizada com a participante em 16 de setembro
        de 2026;
  \item Código-fonte e documentação da versão \textit{Alpha 3i} do aplicativo.
\end{itemize}

\end{document}